% related_work_positioning.tex
% Positioning of K-Bound against concurrent label-free TTA-safety work.
% Drafted 2026-06-19 from primary abstracts (AETTA CVPR'24; Monitoring Risks in TTA,
% Schirmer et al. 2025; Suitability Filter 2025) + ATC (Garg et al. ICLR'22) and
% agreement-on-the-line (Baek et al. NeurIPS'22).
% \input this where the discussion of related work belongs. References are given as
% footnotes so the file compiles without .bib edits; suggested \citep keys are noted
% in comments for later wiring.

\section{Positioning vs.\ concurrent label-free TTA-safety work}\label{sec:positioning}

K-Bound sits at the intersection of three active, label-free research lines. None of
them, individually, answers the question we pose---\emph{when is the adapt-or-not
decision knowable without target labels, and what guarantee can accompany it?}---but
each is close enough that the distinction must be made precisely. We do \emph{not}
claim to be the first to observe that test-time adaptation (TTA) can hurt, nor the
first to apply anytime-valid sequential statistics in this setting; our contribution
is the \emph{decision-theoretic knowability boundary} (Theorem~\ref{thm:imp} and the
one-bit dichotomy) together with a \emph{pre-commitment} adapt/freeze/abstain rule
carrying a controlled false-adapt rate.

\paragraph{(A) Unsupervised accuracy estimation.}
ATC\footnote{Garg et al., \emph{Leveraging unlabeled data to predict out-of-distribution
performance}, ICLR 2022.} and agreement-on-the-line\footnote{Baek et al., \emph{Agreement-on-the-line:
predicting the performance of neural networks under distribution shift}, NeurIPS
2022.}---and more recent gradient- and dispersity-based
estimators---predict a model's accuracy on unlabeled shifted data from heuristic
signals (thresholded confidence, cross-model agreement). They estimate a
\emph{scalar} (accuracy), rely on covariate-shift--type regularity, and are known to
degrade under concept shift. K-Bound differs on two axes. First, our
Theorem~\ref{thm:imp} supplies the \emph{decision-theoretic reason} these estimators
are fundamentally limited: by a Le Cam two-point construction, the sign of the
adaptation benefit is \emph{undecidable} from label-free observations without exactly
one bit of structural supplement. Second, K-Bound targets a \emph{decision}
(adapt/freeze/abstain), not an accuracy number; the estimator is a means, and the
certificate, not the point estimate, is what is reported.

\paragraph{(B) Label-free accuracy estimation \emph{for} TTA / failure detection.}
AETTA\footnote{Lee et al., \emph{AETTA: Label-Free Accuracy Estimation for Test-Time
Adaptation}, CVPR 2024; arXiv:2404.01351.} estimates a TTA model's accuracy via
prediction disagreement between the target model and its dropout inferences, and uses
that estimate to trigger model recovery. This is a heuristic point estimate with no
coverage guarantee. K-Bound differs in kind: it issues a \emph{certified} decision
with a controlled false-adapt rate, and---via the knowability boundary---characterizes
exactly when \emph{any} such label-free estimate can be trusted (it cannot, in
general, under concept shift). We include AETTA-style disagreement as a
\emph{surrogate} decision baseline; the paper is explicit that the true dropout-pass
AETTA and agreement-on-the-line are not computable from our logged statistics, so the
corresponding column is an honest surrogate, not an apples-to-apples re-run.

\paragraph{(C) Anytime-valid risk monitoring for TTA (the nearest neighbor).}
``Monitoring Risks in Test-Time Adaptation''\footnote{Schirmer, Jazbec, Naesseth,
Nalisnick, \emph{Monitoring Risks in Test-Time Adaptation}, arXiv:2507.08721, 2025.}
extends sequential testing with \emph{confidence sequences} to the TTA setting---model
updated at test time, no labels---raising an alarm when a predefined performance
criterion is violated, i.e.\ detecting the point at which a degraded model should be
taken offline and retrained. This shares our anytime-valid machinery (a confidence
sequence is the interval dual of the e-process underlying
Theorem~\ref{thm:cert}), so we state the distinction head-on:
\begin{itemize}
  \item \textbf{Decision object.} Their framework is a \emph{post-hoc} degradation
  alarm---``has the running adapter failed enough to retrain?'' K-Bound is a
  \emph{pre-commitment} trichotomy---``should this adaptation be applied at all; if
  not, freeze or abstain?''---evaluated before the update is trusted.
  \item \textbf{Theoretical core.} Monitoring provides anytime-valid \emph{tracking}
  of a metric (an upper-confidence guarantee). K-Bound additionally provides the
  \emph{impossibility lower bound}---the knowability boundary itself---which no
  monitoring framework supplies and which is the part that explains \emph{why}
  abstention is sometimes the only correct action.
  \item \textbf{Complementarity.} The two compose: K-Bound decides at commit time,
  confidence-sequence monitoring watches afterward. A deployed system could run both.
\end{itemize}
The Suitability Filter\footnote{\emph{Suitability Filter: A Statistical Framework for
Classifier Evaluation in Real-World Deployment Settings}, arXiv:2505.22356, 2025.}
is adjacent: it gates a \emph{fixed} model by hypothesis-testing covariate-shift
``suitability signals'' against a labeled reference, certifying that accuracy has not
dropped beyond a margin. It does not address \emph{adaptation}, nor the concept-shift
impossibility; K-Bound's decision is about whether to \emph{change} the model and is
accompanied by the undecidability characterization.

\paragraph{Summary of the delta.}
\begin{center}
\small
\begin{tabular}{lllll}
\hline
\textbf{Line / method} & \textbf{Decision object} & \textbf{Guarantee} & \textbf{Impossibility?} & \textbf{Role} \\
\hline
ATC / agreement-on-line & accuracy scalar & none (empirical) & not characterized & monitor (static) \\
AETTA (CVPR'24) & accuracy for TTA & none (heuristic) & mitigated heuristically & monitor (post-hoc) \\
Monitoring Risks (2025) & degradation alarm & anytime-valid \emph{upper} & detects, no lower bnd & monitor (post-hoc) \\
Suitability Filter (2025) & fixed-model gate & cov.-shift hyp.\ test & covariate only & monitor (static) \\
\textbf{K-Bound (ours)} & \textbf{adapt/freeze/abstain} & \textbf{false-adapt $\le\alpha$ + lower bnd} & \textbf{characterized (1 bit)} & \textbf{decide (pre-commit)} \\
\hline
\end{tabular}
\end{center}

\paragraph{What is, and is not, novel here (honest scope).}
The anytime-valid certificate machinery is \emph{not} our novel contribution:
confidence-sequence risk monitoring for TTA already exists\footnotemark[\value{footnote}]
and unsupervised accuracy estimation already states the
no-assumption impossibility informally. K-Bound's defensible core is (i) the
\emph{formal} knowability boundary---a Le Cam impossibility tied to the
sign-of-benefit decision, with the exact one-bit supplement that closes it---and (ii)
the \emph{pre-commitment} adapt/freeze/abstain rule with a controlled false-adapt
rate and an honest map of where reliability/drift routing helps. The paper should
therefore lead with the boundary, present the certificate as the operational
instantiation, and cite (A)--(C) as the cluster K-Bound organizes rather than
competes with on machinery.
